\documentclass[10pt,a4paper]{amsart}
\usepackage{amsmath,amssymb,amsthm,graphicx}
% For commutative diagrams we use Xy-pic
\usepackage[all]{xy}

% We use 2cell for 2-commutative diagrams.
\xyoption{2cell}
% Theorem environments.
%
\theoremstyle{plain}
\newtheorem{theorem}[subsection]{Theorem}
\newtheorem{proposition}[subsection]{Proposition}
\newtheorem{lemma}[subsection]{Lemma}

\theoremstyle{definition}
\newtheorem{definition}[subsection]{Definition}
\newtheorem{example}[subsection]{Example}
\newtheorem{exercise}[subsection]{Exercise}
\newtheorem{situation}[subsection]{Situation}

\theoremstyle{remark}
\newtheorem{remark}[subsection]{Remark}
\newtheorem{remarks}[subsection]{Remarks}

\numberwithin{equation}{subsection}

% Macros
%
\def\lim{\mathop{\mathrm{lim}}\nolimits}
\def\colim{\mathop{\mathrm{colim}}\nolimits}
\def\Spec{\mathop{\mathrm{Spec}}}
\def\Hom{\mathop{\mathrm{Hom}}\nolimits}
\def\Ext{\mathop{\mathrm{Ext}}\nolimits}
\def\SheafHom{\mathop{\mathcal{H}\!\mathit{om}}\nolimits}
\def\SheafExt{\mathop{\mathcal{E}\!\mathit{xt}}\nolimits}
\def\Sch{\mathit{Sch}}
\def\Mor{\mathop{\mathrm{Mor}}\nolimits}
\def\Ob{\mathop{\mathrm{Ob}}\nolimits}
\def\Sh{\mathop{\mathit{Sh}}\nolimits}
\def\NL{\mathop{N\!L}\nolimits}
\def\CH{\mathop{\mathrm{CH}}\nolimits}
\def\proetale{{pro\text{-}\acute{e}tale}}
\def\etale{{\acute{e}tale}}
\def\QCoh{\mathit{QCoh}}
\def\Ker{\mathop{\mathrm{Ker}}}
\def\Im{\mathop{\mathrm{Im}}}
\def\Coker{\mathop{\mathrm{Coker}}}
\def\Coim{\mathop{\mathrm{Coim}}}

% Boxtimes
%
\DeclareMathSymbol{\boxtimes}{\mathbin}{AMSa}{"02}

%
% Macros for moduli stacks/spaces
%
\def\QCohstack{\mathcal{QC}\!\mathit{oh}}
\def\Cohstack{\mathcal{C}\!\mathit{oh}}
\def\Spacesstack{\mathcal{S}\!\mathit{paces}}
\def\Quotfunctor{\mathrm{Quot}}
\def\Hilbfunctor{\mathrm{Hilb}}
\def\Curvesstack{\mathcal{C}\!\mathit{urves}}
\def\Polarizedstack{\mathcal{P}\!\mathit{olarized}}
\def\Complexesstack{\mathcal{C}\!\mathit{omplexes}}
% \Pic is the operator that assigns to X its picard group, usage \Pic(X)
% \Picardstack_{X/B} denotes the Picard stack of X over B
% \Picardfunctor_{X/B} denotes the Picard functor of X over B
\def\Pic{\mathop{\mathrm{Pic}}\nolimits}
\def\Picardstack{\mathcal{P}\!\mathit{ic}}
\def\Picardfunctor{\mathrm{Pic}}
\def\Deformationcategory{\mathcal{D}\!\mathit{ef}}
\usepackage{cleveref}
\crefname{originalsection}{original section}{original sections}
\Crefname{originalsection}{Original section}{Original sections}
\makeatletter
\newlabel{section-big-small}{{03X7}{source}{}{sourcepointer-1}{}}
\newlabel{section-big-small@cref}{{[originalsection][0][]03X7}{[1][0][]source}{}{}{}}
\makeatother
\crefname{originallemma}{original lemma}{original lemmas}
\Crefname{originallemma}{Original lemma}{Original lemmas}
\makeatletter
\newlabel{sites-cohomology-lemma-h1-invertible}{{040E}{source}{}{sourcepointer-2}{}}
\newlabel{sites-cohomology-lemma-h1-invertible@cref}{{[originallemma][0][]040E}{[1][0][]source}{}{}{}}
\makeatother
\crefname{originaltheorem}{original theorem}{original theorems}
\Crefname{originaltheorem}{Original theorem}{Original theorems}
\makeatletter
\newlabel{theorem-quasi-coherent}{{03OJ}{source}{}{sourcepointer-3}{}}
\newlabel{theorem-quasi-coherent@cref}{{[originaltheorem][0][]03OJ}{[1][0][]source}{}{}{}}
\makeatother
\crefname{originalproposition}{original proposition}{original propositions}
\Crefname{originalproposition}{Original proposition}{Original propositions}
\makeatletter
\newlabel{descent-proposition-equivalence-quasi-coherent}{{03DX}{source}{}{sourcepointer-4}{}}
\newlabel{descent-proposition-equivalence-quasi-coherent@cref}{{[originalproposition][0][]03DX}{[1][0][]source}{}{}{}}
\makeatother
\title{The first cohomology of the multiplicative group is the Picard group for all the big topologies}
\author{The Stacks Project Authors}
\date{stacks-project:etale-cohomology.tex:local-snapshot-20261011; GNU FDL 1.2 or later}
\begin{document}
\maketitle
\noindent\textit{Editorial note.} All delivered mathematical text is the author's own, copied from the pinned Stacks Project chapter file etale-cohomology.tex; no statement or proof step was rewritten, reordered or removed. The author writes H\_etale\textasciicircum{}1 with the chapter's own \textbackslash{}etale macro and the bare subscript 'syntomic'; both are reproduced verbatim.\par\medskip
\noindent\textit{Editorial note.} The preamble is the chapter's own native preamble excerpt used for this source: the Xy-pic setup at lines 14-18 and the theorem declarations with the mathematical macros at lines 154-217 of sources/stacks-etale-cohomology/preamble.tex. The quoted class line and the xr-hyper externaldocument block are not included because the assembler supplies a plain amsart wrapper; no macro was invented or redefined.\par\medskip
\noindent\textit{Editorial note.} The same-chapter pointer is located at its \textbackslash{}label line inside the pinned chapter file; the three cross-chapter pointers are located in the project's own tag list staged in the snapshot (sources/stacks-chow/tags.txt), at the verbatim lines carrying their canonical tags.\par\medskip
\noindent\textit{Allowed standard prerequisites (editorial).}\begin{itemize}
\item The big and small sites of a scheme and their structure sheaves. Same chapter, Section labelled section-big-small (Stacks tag 03X7); the proof lets tau run over the topologies considered there. Supplies the sites (Sch/X)\_tau and their structure sheaves O\_tau for tau in \{fppf, syntomic, smooth, etale, Zariski\}; the chain of identifications is the counted content.
\item The structure sheaf of a site computes the given sheaf of functions and Picard groups of the topos. Cohomology on Sites, Lemma labelled sites-cohomology-lemma-h1-invertible (canonical Stacks tag 040E); invoked at proof lines 2513-2515. Supplies H\textasciicircum{}1\_tau(X, G\_m) = H\textasciicircum{}1\_tau(X, O\_tau\textasciicircum{}*) = Pic(O\_tau) for each topology tau.
\item Invertible modules are quasi-coherent, and quasi-coherent modules on the big site do not depend on the topology. Same chapter, Theorem labelled theorem-quasi-coherent (source line 1416, Stacks tag 03OJ) and Descent, Proposition labelled descent-proposition-equivalence-quasi-coherent (Stacks tag 03DX); invoked at proof lines 2522-2524. Supply that an invertible O\_tau-module is quasi-coherent and that Pic(O\_tau) = Pic(X); the organisation of the chain is retained in the counted span.
\end{itemize}\medskip
\ifdefined\hypertarget\hypertarget{sourcepointer-1}{}\fi
\noindent\textit{Source reference.} Original section 03X7: sources/stacks-chow/tags.txt, source lines 5138--5138. Cross-chapter pointer to the section of the same chapter that fixes the list of topologies over which tau runs at proof line 2512; the uniform argument for all topologies is retained verbatim in the counted span.\par\medskip
\ifdefined\hypertarget\hypertarget{sourcepointer-2}{}\fi
\noindent\textit{Source reference.} Original lemma 040E: sources/stacks-chow/tags.txt, source lines 5253--5253. Cross-chapter standard prerequisite invoked at proof lines 2513-2515 for H\textasciicircum{}1\_tau(X, G\_m) = H\textasciicircum{}1\_tau(X, O\_tau\textasciicircum{}*) = Pic(O\_tau); the uniformity over tau and the subsequent comparison with Pic(X) are retained in the counted span.\par\medskip
\ifdefined\hypertarget\hypertarget{sourcepointer-3}{}\fi
\noindent\textit{Source reference.} Original theorem 03OJ: sources/stacks-etale-cohomology/pinned-original-source.tex, source lines 1416--1416. Same-chapter pointer offered at proof lines 2523-2524 as one of two ways to see that Pic(O\_tau) = Pic(X); the other route (the descent proposition) is cited in the same sentence and the organisation of the chain is retained in the counted span.\par\medskip
\ifdefined\hypertarget\hypertarget{sourcepointer-4}{}\fi
\noindent\textit{Source reference.} Original proposition 03DX: sources/stacks-chow/tags.txt, source lines 4436--4436. Cross-chapter standard prerequisite invoked at proof lines 2523-2525 as the more precise statement that quasi-coherent modules on the big site are the same for every topology; it supplies that comparison only.\par\medskip
\begin{theorem}
\label{theorem-picard-group}
For any scheme $X$ we have canonical identifications
\begin{align*}
H_{fppf}^1(X, \mathbf{G}_m) & = H^1_{syntomic}(X, \mathbf{G}_m) \\
& = H^1_{smooth}(X, \mathbf{G}_m) \\
& = H_\etale^1(X, \mathbf{G}_m) \\
& = H^1_{Zar}(X, \mathbf{G}_m) \\
& = \Pic(X) \\
& = H^1(X, \mathcal{O}_X^*)
\end{align*}
\end{theorem}

\begin{proof}
Let $\tau$ be one of the topologies considered in
Section \ref{section-big-small}.
By
Cohomology on Sites, Lemma
\ref{sites-cohomology-lemma-h1-invertible}
we see that
$H^1_\tau(X, \mathbf{G}_m) =
H^1_\tau(X, \mathcal{O}_\tau^*) =
\Pic(\mathcal{O}_\tau)$
where $\mathcal{O}_\tau$ is the structure sheaf of the site
$(\Sch/X)_\tau$. Now an invertible $\mathcal{O}_\tau$-module
is a quasi-coherent $\mathcal{O}_\tau$-module.
By Theorem \ref{theorem-quasi-coherent} or the more precise
Descent, Proposition \ref{descent-proposition-equivalence-quasi-coherent}
we see that $\Pic(\mathcal{O}_\tau) = \Pic(X)$.
The last equality is proved in the same way.
\end{proof}

\end{document}
